
% ----------------------------------------------
\chapter{Manipulation}


	 Mesures expérimentales
	
	\begin{lstlisting} language=python, numbers=none]
		
		
		Vol=[4,   3.8,  3.6,  3.4,  3.2,   3,   2.8,  2.6,   2.4, 2.2, 2 ,   1.8,  1.6,  1.4,   1.2,   1, 0.8, 0.6,0.4]
		
		P=  [9.7, 10.2, 10.7,11.2, 11.9, 12.6, 13.3, 14.1, 15.15, 16, 17.1, 18.3 , 19.6, 19.8, 20.1, 19.8, 20, 20 ,20.2]
		P= [i+1 for i in P]
		P1=  [9.8, 10.2, 10.7,11.3, 11.9, 12.6, 13.4, 14.2, 15.1, 16.2, 17.3, 18.5 , 19.9, 20.6, 20.6, 20.7, 20.8, 20.9,20.1]
		P1= [i+1 for i in P1]
		P2 =[9.7, 10.2, 10.7,11.3, 11.9, 12.6, 13.3, 14.1, 15, 16.1, 17.2, 18.4 , 19.5, 20.1, 20.2, 20.3, 20.3, 20.4,20.8]
		P2= [i+1 for i in P2]
		P3 = [9.7, 10.1, 10.7,11.2, 11.8, 12.5, 13.3, 14.1, 15, 16, 17.1, 18.3 , 19.6, 19.9, 20, 20, 20.1, 20.2,20.4]
		P3= [i+1 for i in P3]
		
		
	\end{lstlisting}

% ############################################
\section*{Analyse}

	\subsection*{}
	\textcolor{blue}{Reponse 1: }
	
		Tracer des isothermes d'Andrews sur le meme graphe $ (P,V) $
		
		\begin{figure}[H]
			\centering
			\includegraphics[width=0.5\textwidth]{graphe1.png}
			\caption{Graphe des isthermes d Andrews}
		\end{figure}
		
		Le palier horizontal de liquéfaction correspond à la zone du graphe où la pression reste quasi constante alors que le volume diminue. C’est typiquement la partie des courbes $ P(V) $ pour un volume compris entre $0.6 cm^3 $ et  $ 1.2 cm^3 $
		
		
		
		
		
		
	\subsection*{}
	\textcolor{blue}{Reponse 2: Les volumes du liquide et de la vapeur saturante}\\
	
		\begin{tabular}{|c|c|c|c|c|}\hline
			Isotherme en$ ^0 C$ & Psat(bar)  & Vl ($cm^3$) & Vg ($cm^3$) & Largeur(Vg-Vl) \\ \hline
			20 & 20.8   & 1.4  &  1.8 & 0.4   \\ \hline
			25 & 21.6   & 1.4  &  1.7 & 0.3   \\ \hline
			35 & 21.2   & 1.4  &  1.6 & 0.2   \\ \hline
			40 & 20.9   & 1.4  &  1.6 & 0.2   \\ \hline
		\end{tabular}
		
		
		
	\subsection*{}
	\textcolor{blue}{Reponse 3: Application de la loi des gaz parfaits pour un point en phase gazeuse }
	
		
		
		On choisit un point sur la branche \textbf{vapeur sèche} de l'isotherme $T = 25^\circ$C :
		
		\begin{itemize}
			\item $P = 12,6 \, \text{bar} = 1,26 \times 10^6 \, \text{Pa}$
			\item $V = 3,0 \, \text{cm}^3 = 3,0 \times 10^{-6} \, \text{m}^3$
			\item $T = 25 + 273.15 = 298.15 \, \text{K}$
		\end{itemize}
		
		La loi des gaz parfaits s'écrit :
		\begin{equation}
			PV = nRT
		\end{equation}
		
		On en déduit le nombre de moles :
		\begin{align}
			n &= \frac{PV}{RT} \\
			n &= 1,525 \times 10^{-3} \, \text{mol}
		\end{align}
		
	
	
	\subsection*{}	
	\textcolor{blue}{Reponse 4: Enthalpie de vaporisation à $T = 22,5^\circ$C}
		
		On utilise la relation de Clapeyron :
		\begin{equation}
			\frac{dP_\text{sat}}{dT} = \frac{\Delta H_\text{vap}}{T (V_g - V_l)}
		\end{equation}
		
		Entre $T_1 = 20^\circ$C et $T_2 = 25^\circ$C (isothermes expérimentales) :
		\[
		\Delta H_\text{vap} \approx \frac{\Delta P_\text{sat}}{\Delta T} \times T_m \times (V_g - V_l)
		\]
		
		\textbf{Données :}
		\begin{itemize}
			\item $\Delta T = 5\,\text{K}$
			\item $T_m = 22,5 + 273 = 295,5\,\text{K}$
			\item $V_g - V_l \approx 0,4\,\text{cm}^3 = 4\times10^{-7}\,\text{m}^3$
			\item $\Delta P_\text{sat} \approx 0,8\,\text{bar} = 8\times10^4\,\text{Pa}$
		\end{itemize}
		
		\textbf{Calcul :}
		\begin{align}
			\Delta H_\text{vap} &= \frac{8\times10^4}{5} \times 295,5 \times 4\times10^{-7} \\
			&= 1,89\times10^3 \,\text{J/mol} = 59\,\text{kJ/mol}
		\end{align}
		
		\textbf{Enthalpie massique :}
		\[
		\Delta h_\text{vap} = \frac{59\times10^3}{200,59} = 294\,\text{kJ/kg}
		\]
		
		
		
		
	\subsection*{}
	\textcolor{blue}{Reponse 5: Rapport $\frac{P_c V_{mc}}{RT_c}$ (SF$_6$)}
		
		\textbf{Données SF$_6$}
		\begin{itemize}
			\item $T_c = 45,5^\circ$C $= 318,65\,\text{K}$
			\item $P_c = 37,6\,\text{bar} = 3,76\times10^6\,\text{Pa}$
			\item $R = 8,314\,\text{J/mol·K}$
		\end{itemize}
		
		\textbf{Théorie Van der Waals}
		\[
		V_{mc} = 3b, \quad b = \frac{RT_c}{8P_c}
		\]
		
		\[
		\frac{P_c V_{mc}}{RT_c} = \frac{P_c \cdot 3b}{RT_c} = \frac{3}{8} = 0,375
		\]
		
		\textbf{Valeur expérimentale}
		\[
		V_{mc} \approx 0,32\,\text{L/mol} = 3,2\times10^{-4}\,\text{m}^3\text{mol}^{-1}
		\]
		\[
		\frac{P_c V_{mc}}{RT_c} = \frac{3,76\times10^6 \times 3,2\times10^{-4}}{8,314 \times 318,65} = 0,37
		\]
		
		\textbf{Comparaison : $0,37$ vs $0,375$ (accord excellent !)}
		
		
		
		
	\subsection*{}	
	\textcolor{blue}{Reponse 6:  Constantes de l'équation de Van der Waals}
		
		\textbf{Expressions théoriques}
		\begin{align}
			a &= \frac{27 R^2 T_c^2}{64 P_c} \\
			b &= \frac{R T_c}{8 P_c}
		\end{align}
		
		\textbf{Calcul pour SF$_6$}
		\begin{align}
			b &= \frac{8,314 \times 318,65}{8 \times 3,76\times10^6} = 8,82\times10^{-5} \, \text{m}^3\text{mol}^{-1} \\
			a &= \frac{27 \times (8,314)^2 \times (318,65)^2}{64 \times 3,76\times10^6} = 7,15 \, \text{Pa·m}^6\text{mol}^{-2}
		\end{align}
		
		\textbf{Unités SI}
		\begin{itemize}
			\item $a$ : $\text{Pa·m}^6\text{mol}^{-2}$ (ou $\text{N·m}^4\text{mol}^{-2}$)
			\item $b$ : $\text{m}^3\text{mol}^{-1}$
		\end{itemize}
		
		\begin{center}
			\fbox{\begin{minipage}{0.9\textwidth}
					\textbf{Résumé des résultats :}
					
					\begin{itemize}
						\item $\Delta H_\text{vap} = 59\,\text{kJ/mol} = 294\,\text{kJ/kg}$
						\item $\frac{P_c V_{mc}}{RT_c} = 0,37$ (vs $0,375$ théorique)
						\item $a = 7,15\,\text{Pa·m}^6\text{mol}^{-2}$
						\item $b = 8,82\times10^{-5}\,\text{m}^3\text{mol}^{-1}$
					\end{itemize}
			\end{minipage}}
		\end{center}
		
		
		
		
		
		
		
		
		
		
		
		
		
		
		
		
		
		
		
		
		
		
		
		
		
		
		
		
		
		
		
		
		
		
		
		
		
		
		
		
		
		
		
		
		
		
		
		
		
		
		
		
		
		
		
		
		
		
	\end{document}
% ############################################
\section{Lists}

Unordered Lists:
\begin{itemize}
\item This is an unordered list. 
\item Item 2.
\item It has three items.
\end{itemize}

Ordered List:
\begin{enumerate}
\item This is an ordered list.
\item Item 2.
\item It has three items.
\end{enumerate}

Ordered List (alphabetical):
\begin{enumerate}[label=\Alph*.]
\item This is an ordered list.
\item Item 2.
\item It has three items.
\end{enumerate}

% ----------------------------------------------
\chapter{Figures}\label{ch:figures}

% ############################################
\section{Images}\label{sec:images}

\autoref{fig:image1} shows how to display images.

\begin{figure}[H]
	\centering
	
\includegraphics[width=0.5\textwidth]{f1.jpg}
    \caption{Image}

	\label{fig:image1}
\end{figure}

\begin{figure}[H]
	\centering
	\includegraphics[width=0.5\textwidth]{f4.jpg}
    \caption{Image with Source}
    \captionsource{\cite{mus:16}}	
	\label{fig:image2}
\end{figure}

\begin{figure}[H]
	\centering
	\includegraphics[width=0.5\textwidth]{f7.jpg}
    \caption{Image with Source and Link}
    \captionsource[https://example.org]{J. Doe}	
	\label{fig:image3}
\end{figure}

% ############################################

% ----------------------------------------------

\chapter{Conclusion}\label{ch:conclusion}

% ############################################
\section{Mathematics}\label{sec:maths}

\subsection{Definitions}

\begin{define}	
A function $f: X \rightarrow Y$ is injective if and only if for all $x_1, x_2 \in X$, $x_1 \neq x_2$ implies $f(x_1) \neq f(x_2)$.
\end{define}

\subsection{Formulas}

Formulas can be included inline, e.g. $K_E^{CID}$ to $S$ or as an own block.
$$ otp = f_k(i, k), \qquad i = 1, ..., n $$

\pagebreak
% ############################################
\section{Code Listings}\label{sec:code}

Code can be displayed inline, e.g. \code{UPDATE} or \code{factorial(n)}.

Code listings can have numbered lines, captions, syntax highlighting, ... .

\begin{lstlisting}[caption={[Code Without Line Numbers] Code Without Line Numbers \\ \captionsource{\cite{mus:16}}}, label={lst:codeExample}, language=python, numbers=none]
 0  1  2  3  4  5  6  7  8  9 10 11 12 13 14 15 16 17 
 A  B  C  D  E  F  G  H  I  J  K  L  M  N  O  P  Q  R  
\end{lstlisting}

\begin{lstlisting}[caption={Code With Line Numbers}, label={lst:factorial}, language=Java]
public int factorial(int n) {
	if (n == 1) {
		return 1;
	} else {
		return n * factorial(n - 1) ;
	}
}
\end{lstlisting}
